\chapter{Kaplansky's Density Theorem}

This chapter records the next formalization target: Sakai's Theorem 1.9.1.  It is the natural
next step after the comparison of the ultraweak, strong, and Mackey topologies in Section 1.8.
It also fills the principal gap between that development and the results on ultraweakly closed
ideals and the projection lattice from Section 1.10.

Throughout this chapter,
\[
  S_M=\{x\in M\mid \lVert x\rVert\leq 1\}
\]
denotes the \emph{closed unit ball}.  Sakai calls this set the unit sphere; the distinction is made
explicit here because Mathlib uses separate names for spheres and closed balls.

The theorem is stated first at Sakai's full generality.  A later corollary specializes the auxiliary
test space to the specified predual, which is the form most often called Kaplansky's density
theorem.  None of the unfinished nodes below should introduce a second predual or a competing
operator topology.  They should be built from Mathlib's dual-pairing, weak-topology, convexity, and
closed-ball APIs, together with the existing \texttt{Mackey} and \texttt{Ultraweak} constructions.

\section{Dual-pairing and convex-closure infrastructure}

The two completed results in this section are already in Mathlib.  They are exposed as separate
nodes so that the new proof terminates in existing general functional analysis rather than a
repository-specific reproof.

\begin{theorem}
  \label{thm:krein_milman}
  \lean{closure_convexHull_extremePoints}
  \mathlibok
  Let $K$ be a compact convex subset of a Hausdorff locally convex real topological vector space.
  Then $K$ is the closure of the convex hull of its extreme points.
\end{theorem}

\begin{proof}
  \leanok
  This is the Krein--Milman theorem in Mathlib.
\end{proof}

\begin{theorem}
  \label{thm:convex_closure_same_dual}
  \lean{LinearEquiv.image_closure_of_convex'}
  \mathlibok
  Let two locally convex topologies on linearly equivalent complex vector spaces have their
  continuous duals identified by precomposition with the equivalence.  The equivalence commutes
  with closure on every real-convex set.  In particular, compatible weak and Mackey topologies have
  the same closed convex sets.
\end{theorem}

\begin{proof}
  \leanok
  Mathlib proves this by transporting both closures to the weak topology and applying geometric
  Hahn--Banach separation.
\end{proof}

\begin{lemma}
  \label{lem:bounded_weak_topologies_agree}
  \notready
  Let $B:M\times P\to\mathbb C$ be a bounded dual pairing of normed spaces, and let
  $V\subseteq P$ be a norm-dense complex linear subspace.  On every norm-bounded subset of $M$,
  convergence against all elements of $V$ is equivalent to convergence against all elements of
  $P$.  Equivalently, the weak topologies $\sigma(M,V)$ and $\sigma(M,P)$ induce the same topology
  on each norm-closed ball.
\end{lemma}

\begin{proof}
  \notready
  One implication is immediate.  For the other, fix $p\in P$ and approximate it in norm by
  $v\in V$.  Uniform boundedness of the points of the net bounds the evaluation error
  $|B(x,p-v)|$ uniformly, while convergence against $v$ handles the remaining term.  This is a
  general bounded-dual-pairing lemma and should be added at that level, not in the operator-algebra
  namespace.
\end{proof}

\begin{lemma}
  \label{lem:mackey_continuous_of_dual_map}
  \notready
  Let $B:E\times F\to\mathbb C$ and $B':E'\times F'\to\mathbb C$ be separating dual pairings.  If
  complex-linear maps $T:E\to E'$ and $T':F'\to F$ satisfy
  \[
    B'(Tx,y)=B(x,T'y),
  \]
  then $T$ is continuous from the Mackey topology $\tau(E,F)$ to
  $\tau(E',F')$.
\end{lemma}

\begin{proof}
  \notready
  Pull an absolutely convex weakly compact subset of $F'$ back along $T'$.  Its image is again
  absolutely convex and weakly compact, and the displayed identity identifies the corresponding
  Mackey seminorms.  This is the functorial map theorem missing from the current \texttt{Mackey}
  API; the existing \texttt{Mackey.congrLeft} handles only equivalences.
\end{proof}

\section{Sakai-invariant test spaces}

\begin{definition}
  \label{def:sakai_invariant_test_space}
  \notready
  Let $M$ have specified predual $P$.  A \textbf{Sakai-invariant test space} is a norm-dense
  complex linear subspace $V\subseteq P$ which is closed under the dual actions induced by
  involution and fixed left and right multiplication on $M$: for $v\in V$ and $a\in M$, the
  functionals
  \[
    x\longmapsto \overline{v(x^{*})},\qquad
    x\longmapsto v(ax),\qquad
    x\longmapsto v(xa)
  \]
  are again represented by elements of $V$.
\end{definition}

The formal object should be a predicate on a submodule of the specified predual, with the three
induced maps recorded explicitly.  It should not be a typeclass: a single predual may contain many
useful test spaces, and inference should not choose one globally.

\begin{lemma}
  \label{lem:tau_continuity_of_sakai_invariant}
  \notready
  \uses{def:sakai_invariant_test_space,lem:mackey_continuous_of_dual_map}
  If $V$ is a Sakai-invariant test space, then involution and the maps
  $x\mapsto ax$ and $x\mapsto xa$ are continuous for the Mackey topology $\tau(M,V)$.
\end{lemma}

\begin{proof}
  \notready
  Apply Lemma \ref{lem:mackey_continuous_of_dual_map} to each operation and the corresponding
  induced map on $V$.  Complex conjugation means that involution is naturally real-linear; its
  continuity should be stated through the existing star/topological-star API rather than forced
  into a complex-linear map.
\end{proof}

\section{The contraction transform}

For $a\in M$, define Sakai's contraction transform by
\[
  \kappa(a)=2a(1+a^{*}a)^{-1}.
\]
The next three nodes isolate the $C^{*}$--algebraic part of the argument from the later topology.
This separation is important: the transform and its estimates do not depend on a predual.

\begin{definition}
  \label{def:kaplansky_transform}
  \notready
  For an element $a$ of a unital $C^{*}$--algebra, the \textbf{Kaplansky contraction transform} is
  $\kappa(a)=2a(1+a^{*}a)^{-1}$.
\end{definition}

The public construction may be packaged as a function once invertibility of $1+a^{*}a$ is supplied
canonically by positivity.  Its membership theorem below must also cover a non-unital
$C^{*}$--subalgebra of a unital ambient algebra.

\begin{proposition}
  \label{prop:kaplansky_transform_mem_closedBall}
  \notready
  \uses{def:kaplansky_transform}
  For every $a$ in a unital $C^{*}$--algebra, $1+a^{*}a$ is invertible and
  $\lVert\kappa(a)\rVert\leq 1$.
\end{proposition}

\begin{proof}
  \notready
  Positivity gives $1\leq 1+a^{*}a$ and hence invertibility.  Continuous functional calculus
  reduces the norm estimate to the scalar inequality
  $2\sqrt t/(1+t)\leq1$ for $t\geq0$.  The implementation should reuse Mathlib's positive-order,
  inverse, square-root, and continuous-functional-calculus APIs.
\end{proof}

\begin{lemma}
  \label{lem:kaplansky_transform_mem_subalgebra}
  \notready
  \uses{def:kaplansky_transform,prop:kaplansky_transform_mem_closedBall}
  Let $A$ be a norm-closed, possibly non-unital star subalgebra of a unital $C^{*}$--algebra $M$.
  If $a\in A$, then $\kappa(a)\in A$.  Consequently, $\kappa$ maps $A$ into $A\cap S_M$ even
  when $1\notin A$.
\end{lemma}

\begin{proof}
  \notready
  Express the transform through the non-unital continuous functional calculus, using a scalar
  function which vanishes at the distinguished point.  This avoids adjoining an identity to $A$
  and keeps the result compatible with Mathlib's non-unital $C^{*}$--subalgebra API.
\end{proof}

\begin{lemma}
  \label{lem:kaplansky_transform_fixed_on_extreme}
  \notready
  \uses{def:kaplansky_transform,lem:extreme_partial_isom}
  If $u$ is an extreme point of $S_M$, then $\kappa(u)=u$.
\end{lemma}

\begin{proof}
  \notready
  Lemma \ref{lem:extreme_partial_isom} makes $p=u^{*}u$ a projection.  Since
  $(1+p)^{-1}=1-\tfrac12p$ and $up=u$, direct calculation gives
  $2u(1+p)^{-1}=u$.
\end{proof}

\section{The density theorem}

\begin{proposition}
  \label{prop:kaplansky_transform_in_ball_closure}
  \notready
  \uses{def:sakai_invariant_test_space,lem:bounded_weak_topologies_agree,thm:convex_closure_same_dual,lem:tau_continuity_of_sakai_invariant,prop:kaplansky_transform_mem_closedBall,lem:kaplansky_transform_mem_subalgebra}
  Let $V\subseteq P$ be a Sakai-invariant test space, and let $A$ be a norm-closed star
  subalgebra of $M$ which is dense for $\sigma(M,V)$.  Then for every $a\in M$,
  \[
    \kappa(a)\in\overline{A\cap S_M}^{\,\sigma(M,P)}.
  \]
\end{proposition}

\begin{proof}
  \notready
  Because $A$ is a linear subspace, it is convex.  The weak and Mackey topologies associated to
  the same dual pairing have the same convex closures, so choose a net $(a_i)$ in $A$ converging to
  $a$ for $\tau(M,V)$.  Sakai's resolvent calculation, using the continuity from Lemma
  \ref{lem:tau_continuity_of_sakai_invariant}, shows that
  $\kappa(a_i)\to\kappa(a)$ in $\sigma(M,V)$.  Each $\kappa(a_i)$ belongs to $A\cap S_M$.
  Lemma \ref{lem:bounded_weak_topologies_agree} upgrades this bounded convergence to
  $\sigma(M,P)$--convergence.
\end{proof}

\begin{theorem}
  \label{thm:kaplansky_density_Sak_1_9_1} (Sakai 1.9.1)
  \notready
  \uses{thm:krein_milman,thm:convex_closure_same_dual,prop:kaplansky_transform_in_ball_closure,lem:kaplansky_transform_fixed_on_extreme}
  Let $M$ be a $W^{*}$--algebra with specified predual $P$, let $V\subseteq P$ be a
  Sakai-invariant test space, and let $A$ be a self-adjoint subalgebra of $M$ which is dense for
  $\sigma(M,V)$.  Then $A\cap S_M$ is dense in $S_M$ for the Mackey topology
  $\tau(M,P)$.
\end{theorem}

\begin{proof}
  \notready
  Replacing $A$ by its norm closure preserves the hypotheses and the desired conclusion.
  Proposition \ref{prop:kaplansky_transform_in_ball_closure} and Lemma
  \ref{lem:kaplansky_transform_fixed_on_extreme} show that the ultraweak closure of
  $A\cap S_M$ contains every extreme point of $S_M$.  That closure is convex.  The closed unit ball
  is ultraweakly compact by Banach--Alaoglu, so Theorem \ref{thm:krein_milman} says it is the
  closure of the convex hull of its extreme points.  Hence $A\cap S_M$ is ultraweakly dense in
  $S_M$.  Finally, Theorem \ref{thm:convex_closure_same_dual} identifies the closure of this convex
  set for the compatible ultraweak and Mackey topologies.
\end{proof}

\begin{corollary}
  \label{cor:kaplansky_density_ultraweak}
  \notready
  \uses{thm:kaplansky_density_Sak_1_9_1,thm:sigma_le_s_le_tau}
  Let $A$ be a self-adjoint subalgebra which is ultraweakly dense in a $W^{*}$--algebra $M$.
  Then $A\cap S_M$ is dense in $S_M$ for the Mackey topology.  Consequently it is also dense for
  the strong and ultraweak topologies.
\end{corollary}

\begin{proof}
  \notready
  Take $V=P$ in Theorem \ref{thm:kaplansky_density_Sak_1_9_1}.  The identity subspace is
  automatically Sakai-invariant by ultraweak continuity of involution and fixed multiplication.
  The final implications follow from
  $\sigma(M,P)\leq s(M,P)\leq\tau(M,P)$.
\end{proof}
